\title{Parameter-Efficient Orthogonal Finetuning via Factorization}

\begin{abstract}
The widespread adoption of large foundation models has made training them from the ground up an impractical endeavor due to the immense computational cost. Consequently, developing methods to adapt these robust models to new tasks with greater efficiency has become increasingly vital. This work explores Orthogonal Finetuning (OFT), a systematic approach for adapting models to downstream tasks. While OFT exhibits strong transferability, it remains parameter-intensive, primarily because orthogonal matrices inherently possess high dimensionality. To improve parameter-efficiency, we analyze OFT through the framework of information transfer and distill a set of criteria necessary for reducing parameter usage. Motivated by the efficiency of the Cooley-Tukey fast Fourier transform in communicating information, we introduce an orthogonal parameterization grounded in butterfly structures. Integrating this parameterization with OFT, we develop Orthogonal Butterfly (BOFT), a new finetuning method that achieves high parameter-efficiency. Notably, OFT can be seen as a special instance of BOFT, allowing BOFT to serve as a broader orthogonal finetuning scheme. We systematically evaluate the effectiveness of this method by adapting large-scale vision transformers, prominent language models, and state-of-the-art text-to-image diffusion models across a spectrum of vision and language downstream tasks.
\end{abstract}

\section{Introduction}

The remarkable performance of contemporary models, such as ChatGPT~\cite{brown2020language,chen2021evaluating} and Stable Diffusion~\cite{rombach2022high}, exemplifies the impressive generalization capabilities of large foundation models. The leap in performance of these models is largely driven by a substantial escalation in parameter count (\emph{e.g.}, GPT-3 encompasses approximately 175B parameters). As a consequence, constructing a foundation model from scratch has become increasingly unfeasible for most research groups. The progression of the field, therefore, hinges on the ability to adapt existing foundation models efficiently, bypassing the need for retraining. To facilitate this, models must be tailored to new downstream tasks with resource-conscious adaptation strategies. There are three principal strategies for effective model adaptation: \emph{model finetuning}~\cite{hulora2022,zaken2022bitfit,guo2020parameter,raffel2020exploring,qiu2023controlling,zhang2022adaptive,chavan2023one}, which maintains the original model structure while adjusting a subset of its parameters; \emph{adapter tuning}~\cite{rebuffi2017learning,houlsby2019parameter,pfeiffer2020adapterfusion,lin2020exploring,he2021towards}, where small, trainable modules are appended to the network; and \emph{prompt tuning}~\cite{li2021prefix,lester2021power}, which augments input sequences with tunable prefix tokens. Of these techniques, model finetuning stands out for its straightforwardness and effectiveness, with the added advantage of not increasing inference time.

The main idea underlying model finetuning is to maintain a close resemblance between the pretrained and finetuned models according to specific metrics, thereby retaining the acquired knowledge from pretraining. Traditional finetuning techniques often utilize a small learning rate, a heuristic that aims to limit the divergence between the two model states. Recognizing the inherent structure within weight matrices, more systematic strategies aim to preserve their relational properties—such as the pairwise angular relationships among neurons. Building on this observation, Orthogonal Finetuning~(OFT)~\cite{qiu2023controlling} was proposed as a new framework in which an identical orthogonal transformation is applied across all neurons in a given layer, thus guaranteeing the invariance of pairwise angles during finetuning. While OFT has shown strong results in terms of generalization and convergence for text-to-image diffusion model finetuning~\cite{qiu2023controlling}, it faces the challenge that orthogonal matrices, especially in high dimensions, involve a substantial number of trainable parameters. To mitigate this, OFT employs a block-diagonal construction to reduce parameter count. Nonetheless, this gain in parameter efficiency is accompanied by limitations: the imposed block-diagonal structure fixes the sparsity pattern, and orthogonal transformations are independently executed within each block. As a consequence, such sparsity introduces certain inductive biases—for instance, block-diagonal orthogonal matrices are less expressive and are unable to approximate general linear mappings. Additionally, there remains a fundamental issue: determining an optimal sparsity structure for these matrices is still unresolved.

The central challenge lies in constructing a dense orthogonal matrix that maintains parameter efficiency. Although it initially appears unattainable because representing a dense orthogonal matrix in $d$ dimensions typically involves $\mathcal{O}(d^2)$ parameters, we introduce a novel strategy: synthesizing a dense orthogonal matrix through the product of several sparse, factorized matrices. This method is motivated by the idea that computation can substitute for storage, thereby reducing the total number of learnable parameters. Because our representation expresses the orthogonal matrix as the sequential multiplication of sparse factors, these multiplications must be executed during every training cycle. To contextualize matrix factorization, we conceptualize the formation of a dense orthogonal matrix as a process akin to information transmission. Framed this way, building a dense orthogonal matrix by composing sparse matrices corresponds to the challenge of transmitting information across a grid-like graph. This viewpoint on information transmission enables us to design a range of sparse matrix factorizations that control parameter count while preserving the capacity to realize dense orthogonal matrices.

To optimize parameter usage within our information transmission perspective, we draw on butterfly network architectures found in the Cooley-Tukey fast Fourier transform algorithm~\cite{cooley1965algorithm}, which facilitates efficient data exchange between distinct nodes~\cite{klasing1994broadcasting}. Notably, the butterfly structure inherent in the Cooley-Tukey method offers a natural approach for sparse matrix factorization, often termed butterfly factorization. Employing this technique, it becomes possible to produce a $d$-dimensional dense matrix as the product of $\mathcal{O}(\log d)$ sparse matrices, each containing only $\mathcal{O}(d)$ nonzero components. By ensuring each sparse factor is itself orthogonal, this technique yields a highly efficient orthogonal parameterization, requiring only $\mathcal{O}(d \log d)$ parameters—a dramatic reduction from the full parameter count. Harnessing this compact parameterization, we introduce Orthogonal Butterfly~(BOFT), a novel finetuning approach characterized by parameter efficiency. BOFT generalizes OFT, which emerges as a specific case within this broader framework, offering a versatile orthogonal finetuning methodology. Both the block-diagonal and butterfly arrangements exploit sparsity, incorporating it into the orthogonal matrices to diminish the number of parameters that must be optimized. In contrast to low-rank matrices as utilized in LoRA~\cite{hulora2022}, our approach is anchored in sparse matrix structures; both low-rank and sparse representations are established structured matrix classes~\cite{candes2011robust} for achieving parameter efficiency.

In contrast to OFT, which leverages a block-diagonal structure to mediate between model expressivity and regularity, BOFT employs the butterfly architecture, offering a more seamless interpolation between elements of the full orthogonal group (for expressivity) and the identity matrix (for regularization). This approach facilitates the identification of a more compact hypothesis space within the orthogonal group for subsequent learning tasks. Given that butterfly structures underpin numerous efficient linear operators, including discrete Fourier and discrete cosine transforms, we contend that adopting this structured paradigm for parameter efficiency implicitly encodes beneficial inductive biases that can enhance model generalization and help mitigate overfitting. This reasoning stems from the capability of butterfly structures to efficiently approximate a variety of classical linear transforms—something block-diagonal structures, as used in OFT, inherently cannot achieve. Our key contributions can be summarized as follows:

\begin{itemize}[leftmargin=*,nosep]
    \item We address the challenge of parameter efficiency in orthogonal finetuning by introducing a new information transmission perspective, which reformulates the design of a dense, parameter-efficient orthogonal matrix as an information flow problem on a grid graph.
    \item Drawing inspiration from the butterfly structures characteristic of the Cooley-Tukey algorithm, we introduce Orthogonal Butterfly—an orthogonal finetuning approach optimized for parameter efficiency—within our information transmission framework.
    \item We present several theoretical results elucidating how BOFT manages to drastically lower the number of tunable parameters, yet maintains robust expressivity and training stability. Furthermore, the underlying matrix decomposition in BOFT yields an interesting form of weight interpolation (see Figure~\ref{fig:boft_interpolation}).
    \item To the best of our knowledge, we are the first to extend orthogonal finetuning beyond the domain of controllable text-to-image synthesis~\cite{qiu2023controlling}, demonstrating its effectiveness as a versatile technique for model adaptation. We evaluate BOFT across a variety of adaptation scenarios in computer vision and natural language processing, consistently surpassing contemporary parameter-efficient baselines and affirming its advantages in generalization and efficiency.
\end{itemize}

\section{Related Work}

\textbf{Parameter-efficient finetuning (PEFT)}. As foundation models (\emph{e.g.}, \cite{brown2020language,radford2021learning,rombach2022high,kirillov2023segment}) continue to scale in size and capabilities, the challenge of adapting these models to specific downstream tasks with minimal trainable parameters has garnered substantial attention~\cite{treviso2023efficient,ding2023parameter,lialin2023scaling}. A variety of PEFT strategies have been explored~\cite{houlsby2019parameter,aghajanyan2020intrinsic,hulora2022,edalati2022krona,wang2022adamix,gheini2021cross,zaken2022bitfit,guo2020parameter,sung2021training,ansell2022composable,lester2021power,li2021prefix,vu2022spot,he2021towards,mao2021unipelt,karimi2021compacter,liu2022few,sung2022lst,chen2022parameter,jia2022visual,chen2022adaptformer,zhang2022neural,jie2023fact,lian2022scaling,luo2023towards}, with reparameterization-based methods~\cite{aghajanyan2020intrinsic,hulora2022,edalati2022krona} particularly pertinent to our approach. LoRA~\cite{hulora2022} modifies the pretrained weights by introducing a sum with the product of two low-rank matrices, yielding strong results for natural language applications. While LoRA assigns a fixed rank to all such matrices, alternative schemes~\cite{zhang2022adaptive,valipour2022dylora,zhang2023increlora} adaptively allocate ranks at each layer, providing better control over parameter usage. The simplicity of low-rank reparameterizations has contributed to their widespread adoption \cite{dettmers2023qlora,zi2023delta,chavan2023one}. Motivated by research linking hyperspherical energy to generalization~\cite{liu2018learning,liu2021orthogonal}, \cite{qiu2023controlling} introduce orthogonal finetuning (OFT), which provides an alternative means of adaptation by learning an orthogonal transformation within each layer. OFT has demonstrated enhanced generalization and stability in training, particularly for text-to-image diffusion models, compared to LoRA. However, OFT often incurs a larger number of trainable parameters relative to LoRA. Consequently, there is value in further optimizing the parameter efficiency of OFT. Furthermore, the extent to which OFT applies beyond text-to-image diffusion tasks~\cite{qiu2023controlling} remains unexplored. BOFT tackles the parameter efficiency concern of OFT through butterfly factorization, thereby enabling a broader demonstration of orthogonal finetuning for a diverse range of adaptation scenarios.

\textbf{Butterfly structures}. The radix-2 Cooley-Tukey algorithm, which decomposes an $N$-point discrete Fourier transform into two subproblems of size $\frac{N}{2}$, naturally yields a butterfly architecture that can be mathematically described as a product of several sparse matrices (collectively forming a butterfly matrix). Such matrices have served various purposes, including parameterizing orthogonal matrices to circumvent pivoting during Gaussian elimination and to boost computational efficiency~\cite{parker1995random}, improving the stability of recurrent neural network training~\cite{jing2017tunable}, and facilitating kernel approximations~\cite{munkhoeva2018quadrature}. Methods from \cite{dao2019learning,dao2019kaleidoscope} have shown how to learn rapid linear transformations using butterfly parameterizations, while \cite{chen2022pixelated,dao2022monarch} applied butterfly matrices for effective neural network training. Butterfly frameworks have also contributed to accelerating matrix-vector operations~\cite{michielssen1996multilevel,o2010algorithm}, supporting data-sparse matrix representations~\cite{li2015butterfly}, and optimizing network transmission protocols~\cite{klasing1994broadcasting,li2003linear,rajasingh2016transmission}. Distinct from this previous literature, our work leverages butterfly structures specifically to improve the parameter efficiency of OFT.

\section{An Information Transmission View on Orthogonal Finetuning}

We begin by outlining key aspects of OFT. In order to adapt a pretrained weight matrix, OFT employs a reparameterization where the updated weight matrix is expressed as a product between a learnable orthogonal matrix and the fixed, pretrained weights. Unlike LoRA, which implements an additive update using a low-rank matrix, OFT updates weights through a multiplicative process involving an orthogonal matrix. LoRA achieves parameter efficiency by constraining updates to low-rank forms, while the standard OFT approach adopts a block-diagonal orthogonal architecture~\cite{qiu2023controlling}; parameter efficiency is enhanced by reducing the block size within the diagonal. Figure~\ref{fig:comp_lora} visually contrasts these two strategies. The primary rationale for employing orthogonal transformations in finetuning is to maintain the angular relationships among neurons~\cite{liu2018learning,liu2021orthogonal,qiu2023controlling}, thereby retaining the pretrained model's semantic knowledge.

Formally, OFT seeks to learn an orthogonal matrix $\bm{R}\in\mathbb{R}^{d\times d}$ for a given pretrained linear layer $\bm{W}^0\in\mathbb{R}^{d\times n}$. The forward computation is thus altered from $\bm{z}=(\bm{W}^0)^\top\bm{x}$ to $\bm{z}=(\bm{R}\bm{W}^0)^\top\bm{x}$, with $\bm{x}\in\mathbb{R}^d$ representing the input and $\bm{z}\in\mathbb{R}^n$ the resulting output. For zero initialization, $\bm{R}$ is set to the identity matrix at the outset. To enforce that $\bm{R}$ remains orthogonal during finetuning, we utilize the Cayley parameterization, following the approach in~\cite{qiu2023controlling,liu2021orthogonal}: $\bm{R}=(\bm{I}+\bm{Q})(\bm{I}-\bm{Q})^{-1}$, where the skew-symmetric condition $\bm{Q}=-\bm{Q}^\top$ holds. Parameter efficiency is realized by modeling $\bm{R}$ as block-diagonal, $\text{diag}(\bm{R}_1,\bm{R}_2,\cdots,\bm{R}_r)$, where each $\bm{R}_i\in\mathcal{R}^{b\times b}$ is a small orthogonal block and $br = d$. This block-diagonal design restricts parameters but imposes an artificial partition on each neuron (i.e., on each column vector of $\bm{W}^0$), grouping its dimensions into $r$ segments and independently transforming each group using a distinct orthogonal matrix. While practical experiments show the utility of such block-diagonal structures, the arbitrary division of neuron dimensions according to their indices lacks theoretical justification, suggesting that employing dense orthogonal matrices could be advantageous. This leads us to a fundamental question: \emph{Is it possible to devise a dense orthogonal matrix that maintains parameter efficiency?}

To explore this issue, we propose expressing a dense orthogonal matrix as the product of several sparse orthogonal matrices. To create a cohesive and accessible framework for analyzing the sparsity structures within orthogonal matrix factorization, we conceptualize the task of constructing a dense orthogonal matrix in OFT as analogous to an information transmission process. Concretely, forming a dense matrix $\bm{R}\in\mathbb{R}^{d\times d}$ by multiplying $m$ square matrices, so that $\thickmuskip=2mu \medmuskip=2mu \bm{R}=\bm{B}_m\bm{B}_{m-1}\cdots\bm{B}_1$, can be mapped to the transmission of information across a lattice comprising $d\times (m+1)$ nodes, depicted in Figure~\ref{fig:info_trans}. This perspective on information transfer arises from observing that a dense $d$-dimensional square matrix can represent fully connected pathways from one set of $d$ nodes to another. In $\bm{R}$, if $R_{ij}$ is zero, then there is no pathway for information from node $j$ to node $i$; if $R_{ij}$ is nonzero, a transmission is possible. Thus, decomposing the dense matrix $\bm{R}$ as $\bm{B}_m\bm{B}_{m-1}\cdots\bm{B}_1$ can be interpreted as a stepwise propagation of information governed by the sequence of graphs defined by each $\bm{B}_i$. The information propagates beginning with $\bm{B}_1$ and sequentially through to $\bm{B}_n$. In the specific example illustrated in Figure~\ref{fig:info_trans}, we consider the decomposition $\bm{R}=\bm{B}_5\bm{B}_4\bm{B}_3\bm{B}_2\bm{B}_1$, and visualize both the corresponding sparsity patterns and the resulting graph structure. The depicted graph results from unfolding the matrix multiplications across layers. In this framework, $\bm{B}_i$ acts as the adjacency matrix representing possible transmissions from nodes at layer $i$ to those at layer $i+1$: specifically, entry $(j_1,j_2)$ in $\bm{B}_i$ indicates whether there is a directed link from node $j_2$ in layer $i$ to node $j_1$ in layer $i+1$ (with zero indicating the absence of a link). For the composition $\bm{B}_5\bm{B}_4\bm{B}_3\bm{B}_2\bm{B}_1$ to be fully dense, information originating at any node in the initial layer must reach all nodes in the sixth layer. If instead we examine the product $\bm{R}=\bm{B}_4\bm{B}_3\bm{B}_2\bm{B}_1$, connecting source nodes at the first layer to target nodes at the fifth, we observe that there is no viable information pathway from node 1 to node 3. Accordingly, the sparsity of $\bm{B}_4\bm{B}_3\bm{B}_2\bm{B}_1$ yields a zero in position $(3,1)$. The same pattern is observed in the sub-products $\bm{R}=\bm{B}_3\bm{B}_2\bm{B}_1$ and $\bm{R}=\bm{B}_2\bm{B}_1$, reflecting the alignment between the induced graphs and the zero structure. In the general case, to guarantee the density of a matrix $\bm{R}\in\mathbb{R}^{d\times d}$, the sequence of $m$ factor matrices $\bm{B}_m,\cdots,\bm{B}_1$ must collectively encode directed edges over a $d\times (m+1)$ network where each edge only links nodes in neighboring layers (i.e., adjacent columns). This structure must enable each starting node in the first layer to communicate with every terminal node in the $(m+1)$-th layer.

The perspective of information transmission offers valuable insights into the patterns of sparsity that emerge within factorization matrices in OFT. As illustrated in Figure~\ref{fig:blk_diag}, the original OFT employs a block-diagonal configuration for $\bm{R}$. While this design helps decrease the overall count of trainable parameters, it inherently limits the construction of $\bm{R}$ to be sparse, not dense. The objective, therefore, is to generate a dense orthogonal matrix by composing $m$ sparse orthogonal matrices, achieving this with a minimal set of effective trainable parameters. Within the framework of information transmission, two principal requirements guide this endeavor: \emph{(i)} \textbf{dense connectivity}, ensuring that each node at the input layer is connected via at least one route to every node at the output layer; and \emph{(ii)} \textbf{minimum free edges}, which dictates minimizing the number of connections permitted by the orthogonality constraint. Orthogonality imposes significant restrictions on edge formation between successive layers. Specifically, for each matrix $\bm{B}_i$ to possess full rank—a precondition for orthogonality—there must be exactly $d$ connections to establish a bijective mapping between the $i$-th layer’s nodes and those of the $(i-1)$-th layer. Consequently, at least $d$ edges are required between any two adjacent layers (see Figure~\ref{fig:info_trans} for the example where $d = 4$). Importantly, these mandatory $d$ edges, essential for orthogonality, are not considered trainable since, in a $d \times d$ orthogonal matrix with $d$ nonzero entries, these entries are constrained to values of $\pm 1$. Orthogonality reduces parameterization compared to unconstrained matrices, which introduces dependencies between the edges. For instance, in the $2 \times 2$ orthogonal case, setting the $(1,1)$ entry to zero obliges the $(2,2)$ entry to be zero as well—meaning the absence of a single edge may enforce the absence of another. Accordingly, the orthogonality condition can alter the feasibility of particular edge configurations, sometimes necessitating the addition or elimination of edges. The information transmission paradigm is particularly pertinent to OFT, as visualizing the sparsity structure (i.e., the non-zero, trainable entries of $\bm{R}$) is critical; only the locations of these elements matter, not their specific numeric values.

A straightforward approach to densely connect two layers involves $\mathcal{O}(d^2)$ edges (essentially, implementing a full dense orthogonal matrix), resulting in $d^2-d$ edges—for instance, when $d=4$, there are 12 such edges. Figure~\ref{fig:info_trans} demonstrates an example where a suitable matrix factorization requires just $10$ edges, a reduction compared to the single dense orthogonal matrix. This methodology allows for analyzing the parameter-efficiency of OFT through a graphical lens, making it convenient to devise efficient factorizations under this framework. We draw on the architecture of butterfly graphs~\cite{cooley1965algorithm}—notably found in the Cooley-Tukey algorithm—which enable all-to-all connections between $d$ sources and $d$ sinks with only $\mathcal{O}(d\log d)$ edges. In contrast, the configuration depicted in Figure~\ref{fig:info_trans} utilizes $10$ edges, but the butterfly network accomplishes complete connectivity with merely $8$ edges. We will next discuss the manner in which butterfly structures can enhance parameter efficiency.

\section{Orthogonal Parameterization by Butterfly Factorization}

The butterfly architecture originates from the Cooley-Tukey algorithm and is essential for efficiently computing the fast Fourier transform. In a Fourier transform, localized frequency domain perturbations yield widespread changes across the spatial domain, mirroring the demands of our information transmission task—in which every node at the initial layer must reach all nodes at the final layer. Additionally, the butterfly configuration has gained prominence in network design for high-throughput communication~\cite{solihin2015fundamentals,li2011linear}. Consider $k\geq 2$ as a power of $2$; the \emph{butterfly factor} $\bm{B}^F(k)$ is then specified as

\begin{equation}\label{eq:but_factor}
\begin{aligned}
    \bm{B}^F(k)=\begin{bmatrix}
    \text{diag}(\bm{d}_1) & \text{diag}(\bm{d}_2) \\
    \text{diag}(\bm{d}_3) & \text{diag}(\bm{d}_4)
    \end{bmatrix}\in\mathbb{R}^{k\times k},
\end{aligned}
\end{equation}

where for each $i$, the vectors $\bm{d}_i\in\mathbb{R}^{\frac{k}{2}}$. Subsequently, for $d=2^N$, the $d$-dimensional \emph{butterfly matrix} $\bm{B}(d)\in\mathbb{R}^{d\times d}$ is recursively built as follows:

\begin{equation}
\begin{aligned}
    \!\!\bm{B}(d)=\tilde{\bm{B}}(d,d)\!\cdot\!\begin{bmatrix}
    \bm{B}_1(\frac{d}{2})\!\!\!\!\! & \bm{0} \\
    \bm{0}\!\!\!\!\! &\bm{B}_2(\frac{d}{2})
    \end{bmatrix}= \tilde{\bm{B}}(d,d)\tilde{\bm{B}}(d,\frac{d}{2})\cdots\tilde{\bm{B}}(d,2),
\end{aligned}
\end{equation}

Here, $\bm{B}_1(\frac{d}{2})$ and $\bm{B}_2(\frac{d}{2})$ denote two butterfly matrices, each of dimension $\frac{d}{2}$. We introduce the concept of a \emph{butterfly component} by defining $\thickmuskip=2mu \medmuskip=2mu \tilde{\bm{B}}(d,k)=\text{diag}(\bm{B}^F_1(k),\cdots,\bm{B}^F_{d/k}(k))$, a $d\times d$ block-diagonal matrix comprising blocks of size $k$, with each $\bm{B}^F_i(k)$ representing a butterfly factor as specified in Equation~\ref{eq:but_factor}. With this foundation, the butterfly matrix can be employed to represent an orthogonal matrix, provided that every multiplicative component $\tilde{\bm{B}}(d,k)$ in the butterfly structure $\bm{B}(d)$ is orthogonal. Focusing initially on the case of block size $2$, $\tilde{\bm{B}}(d,2)$, we observe that orthogonality can be easily enforced by parameterizing each $2\times2$ block via the Cayley transform (or a two-dimensional rotation), as demonstrated in prior works \cite{liu2021orthogonal,qiu2023controlling}. Since the non-zero structure of any butterfly component can be regarded as a permutation of that in $\tilde{\bm{B}}(d,2)$, the parameterization for orthogonality naturally extends to all butterfly components. This method enables a highly efficient parameterization of orthogonal matrices using an assembly of $2\times2$ orthogonal blocks. Building on \cite{chen2022pixelated}, we further generalize the butterfly construction and define a block butterfly component $\tilde{\bm{B}}^b(d,k)$ in which each element of $\bm{d}_i$ is replaced with a $b\times b$ matrix. To ensure the orthogonality of $\tilde{\bm{B}}^b(d,2)$, we parameterize each $2b\times 2b$ block to be orthogonal. The non-zero structure of the generalized butterfly components $\tilde{\bm{B}}^b(d,k)$ for $k>2$ can similarly be viewed as block-wise permutations of that of $\tilde{\bm{B}}^b(d,2)$, making it straightforward to parameterize them as orthogonal as well. Together, these elements constitute the forward computation in BOFT as follows:

\begin{equation}
    \begin{aligned}\nonumber
    \bm{z}=\big(\bm{R}(m,b)\cdot\bm{W}^0\big)^\top\bm{x},~~~~\text{s.t.}~\bigg{\{}\bm{R}(m,b)=\prod_{i=1}^m \tilde{\bm{B}}^b_{(d,i)}~~\&~~\underbrace{\big(\tilde{\bm{B}}^b_{(d,j)}\big)^\top\tilde{\bm{B}}^b_{(d,j)}=\tilde{\bm{B}}^b_{(d,j)}\big(\tilde{\bm{B}}^b_{(d,j)}\big)^\top\!=\bm{I}_d}_{\forall j\in [1, m]}\bigg{\}},
    \end{aligned}
\end{equation}

In this formulation, we use the shorthand $\tilde{\bm{B}}^b(d,2^{m - i+1}) = \tilde{\bm{B}}^b_{(d,i)}$, and $\bm{I}_d$ stands for the $d \times d$ identity matrix. The orthogonal matrix $\bm{R}(m,b)\in\mathbb{R}^{d\times d}$ is constructed as a sequential product of orthogonal butterfly components. For succinctness, BOFT parameterized by $\bm{R}(m,\frac{b}{2})$ is referred to as BOFT($m$,$b$) with $b\geq 2$. When $m=1$, the architecture BOFT($1$,$b$) becomes equivalent to the block-diagonal OFT~\cite{qiu2023controlling}, using block size $b$. For BOFT($1$,$d$), the model reduces to the classic OFT~\cite{qiu2023controlling} with an unconstrained orthogonal matrix of full dimension. The special case BOFT($\log \frac{2d}{b}$,$b$) applies the block butterfly matrix $\bm{B}^{\frac{b}{2}}(d)$ as $\bm{R}$, thereby creating a dense orthogonal matrix. Broadly speaking, BOFT($m$,$b$) involves $\frac{1}{2}(b-1)dm$ effective trainable parameters when fine-tuning a linear layer of dimensions $d\times n$. For the choice $m=\log d, b=2$ (i.e., using butterfly matrices), BOFT achieves parameter efficiency with complexity $\mathcal{O}(d\log d)$. By contrast, the standard OFT relying on full orthogonal matrices requires $\mathcal{O}(d^2)$ parameters, while the block-diagonal OFT with $r$ blocks incurs $\mathcal{O}(bd)$ parameter costs. Therefore, OFT necessitates $b=d$ to attain dense orthogonality, but BOFT is more flexible and can realize dense matrices with arbitrary $b$.

\textbf{Identity initialization for BOFT}. In practice, fine-tuning generally starts from the pretrained model to avoid excessive drift. For example, LoRA initializes the low-rank adaptation weights at zero. In the case of BOFT, all butterfly matrix components are initialized as identity matrices (meaning that the associated skew-symmetric matrices in the Cayley transform start as zero).

\textbf{Multiplicative Dropout for BOFT}. LoRA~\cite{hulora2022} incorporates Dropout for its low-rank adaptation to mitigate overfitting, with conventional Dropout~\cite{srivastava2014dropout} working seamlessly. However, since BOFT updates weights multiplicatively, standard Dropout is not directly applicable. To solve this, we introduce a specialized multiplicative Dropout scheme for BOFT. Given that $\bm{R}(m,b)$ is assembled from $m$ butterfly components, each of which can be rearranged as $2b \times 2b$ block-diagonal orthogonal matrices, our method randomly selects $p_1$ percent of butterfly components and $p_2$ percent of diagonal blocks within each, and sets the selected blocks to identity matrices.

\section{Intriguing Insights and Discussions}

\textbf{Expressivity of BOFT}. The butterfly architecture, when combined with suitable permutations, is capable of exactly replicating several well-known fast linear transforms~\cite{dao2019learning,dao2019kaleidoscope} (such as the fast Fourier transform and Hadamard transform). However, its capacity to approximate arbitrary orthogonal matrices, particularly those that are dense and randomly generated, has not been thoroughly characterized. To probe this, we carry out experiments approximating a randomly sampled dense orthogonal matrix~\cite{anderson1987generation} of size $512\times 512$. Figure~\ref{fig:para_eff} reports the averaged results across 10 independent runs. Here, the y-axis reflects the approximation error, while the x-axis shows the effective number of trainable parameters. Each color-coded curve corresponds to BOFT models employing the same block size, where the leftmost marker indicates the performance of BOFT($1$,$b$)—that is, the conventional block-diagonal OFT of block size $b$. BOFT exhibits increased parameter efficiency compared to OFT; for instance, BOFT($9$,$2$) achieves greater expressiveness than BOFT($1$,$16$) while requiring significantly fewer parameters. In general, configurations with larger $m$ and smaller $b$ tend to use parameters more effectively—BOFT($6$,$4$), for example, achieves error comparable to BOFT($2$,$16$) with substantially reduced parameter counts. The butterfly construction thus encodes a richer set of orthogonal matrices than does a purely block-diagonal configuration, implying that BOFT, given equivalent block size, surpasses OFT in expressive capacity.

\begin{theorem}[Expressivity of BOFT]\label{thm:exp_boft}
    BOFT possesses greater expressive power than OFT when the block size is fixed. Approximating all orthogonal matrices of size $d$ with butterfly matrices can be realized through products of the form $\bm{B}_{d-1,1}(d)\bm{B}^\top_{d-1,2}(d)\cdots\bm{B}_{1,1}(d)\bm{B}^\top_{1,2}(d)$, where each $\bm{B}_{i,j}(d)$ is a butterfly matrix.
\end{theorem}

Theorem~\ref{thm:exp_boft} provides a pathway to a generalized BOFT formulation: the ultimate orthogonal transformation may be constructed as $\thickmuskip=2mu \medmuskip=2mu \bm{R}^G(m_1,b_1,m_2,b_2,l)=\bm{R}_{l,1}(m_1,b_1)\bm{R}_{l,2}^T(m_2,b_2)\cdots\bm{R}_{1,1}(m_1,b_1)\bm{R}_{1,2}^T(m_2,b_2)$, with $\bm{R}_{i,j}^T(m,b)$ denoting an orthogonal factor as specified in BOFT. When parameters are set such that $m_1 = m_2 = \log d$, $b_1 = b_2 = 2$, and $l = d-1$, the resulting composition $\bm{R}^G(m_1,b_1,m_2,b_2,l)$ is capable of representing every orthogonal matrix—a construction that aligns with the so-called kaleidoscope hierarchy~\cite{dao2019kaleidoscope}. Nevertheless, while enhancing expressiveness, this does not invariably translate to improved downstream finetuning outcomes. Indeed, full finetuning—despite its expressive universality—often underperforms. Thus, balancing the model’s expressive capacity with structural regularity is essential to robust finetuning generalization. By endowing the finetuning procedure with structured priors, BOFT expands the set of available tunable parameters, facilitating a more favorable compromise between expressiveness and inductive bias.

\textbf{Spectral properties}. Compared to LoRA, orthogonal finetuning typically achieves superior spectral characteristics, as it exactly maintains the spectral norm of the original pretrained weight matrix $\bm{W}^0$. To illustrate this, consider the singular value decomposition $\bm{W}^0=\bm{U}\bm{\Sigma}\bm{V}^\top$, where $\bm{U}$ and $\bm{V}$ are orthogonal matrices and $\bm{\Sigma}$ is a diagonal matrix containing the singular values. Both OFT and BOFT update the weights by pre-multiplying an orthogonal matrix $\bm{R}$, resulting in new weights $\bm{R}\bm{U}\bm{\Sigma}\bm{V}^\top$; this procedure leaves the leading singular value (\emph{i.e.}, the spectral norm of $\bm{W}^0$) unaffected. Preserving the spectral norm has been demonstrated to promote greater training stability and improved generalization~\cite{miyato2018spectral,yoshida2017spectral}. Additional theoretical insights regarding these properties can be found in Appendix~\ref{app:sn_property}.

\textbf{Orthogonal finetuning as learning bilinear similarity}. In the context of BOFT, the procedure can be expressed as optimizing a bilinear similarity measure $\bm{w}^0_i\bm{R}\bm{x}$, where $\bm{w}^0_i$ corresponds to the $i$-th neuron (\emph{i.e.}, a column of the weight matrix $\bm{W}^0$). Essentially, BOFT is tasked with identifying an orthogonal matrix $\bm{R}$ that defines the bilinear similarity, subject to stringent regularity constraints (\emph{i.e.}, orthogonality). This process is fundamentally linked to distance metric learning~\cite{xing2002distance} as well as to the theory of bilinear forms~\cite{roman2005advanced}.

\textbf{Inductive bias and generalization in BOFT}. Because the transformation $\bm{R}(m,b)$ in BOFT typically belongs to a structured subgroup of the broader orthogonal group, it serves to limit the overall hypothesis space and thus introduces a pronounced inductive bias. We contend that the structured inductive bias conferred by the butterfly factorization has a positive impact on generalization, owing to its structural resemblance to many classical linear transforms~\cite{dao2019kaleidoscope}—including the discrete Fourier, discrete sine/cosine, and Hadamard transforms. Furthermore, the sparse structure arising from the matrix factorization in BOFT can lead to further implicit inductive biases~\cite{gunasekar2017implicit,li2018algorithmic,liu2019neural}.

\textbf{Relation to butterfly-based sparse training}. Several previous studies~\cite{dao2019learning,dao2019kaleidoscope,chen2022pixelated,dao2022monarch} have investigated sparse training approaches utilizing butterfly parameterizations. The predominant approach in these works involves directly reparameterizing weight matrices with the butterfly structure and training models from the initial randomization. In particular, \cite{dao2022monarch} explores a finetuning paradigm where pretrained weights are first mapped onto a variant of butterfly matrices, followed by optimization of the resultant projected elements for downstream applications. In contrast, BOFT introduces a distinct finetuning mechanism that applies transformations to the weights via layer-shared matrix operators.

\input{rebuttal-results}

\section{Concluding Remarks and Limitations}

In this work, we introduce Orthogonal Butterfly, a broadly applicable and parameter-efficient method for finetuning foundation models, constructed upon the butterfly architecture. Our principal observation is that higher parameter-efficiency can be attained by expressing a dense orthogonal transformation as a sequence of sparse orthogonal matrices. To facilitate feasible decompositions of matrices, we develop a graphical framework rooted in information transmission. Within this conceptual setting, we observe that butterfly structures fulfill the requirements for sparse orthogonal matrix decomposition particularly well. We provide comprehensive empirical validation for the effectiveness of BOFT across various tasks including finetuning large language models, vision models, and text-to-image generative models. Experimental results underscore BOFT's robustness as a universal finetuning strategy.

Nonetheless, BOFT is not devoid of shortcomings. Because the resultant orthogonal matrix in BOFT arises from the composition of several orthogonal matrices, the computational cost during training exceeds that of OFT. Addressing this increased training overhead is a direction for future research. Importantly, after finetuning, the orthogonal matrices discovered by BOFT can be fused into the base model without incurring additional inference costs. Additionally, it remains uncertain whether the butterfly network represents the optimal architecture for information transmission. Our proposed information transmission perspective invites further exploration and draws a parallel with computer networking research, where optimal network topologies for information flow are rigorously analyzed. We anticipate that alternative, potentially more efficient, network configurations may be exploited to construct dense orthogonal matrices.